% Databeheer-Database-Design.tex
% =============================================================================
%     SAMENVATTING DATABEHEER: ERD & DATABASE DESIGN -- KU Leuven
%     Gebaseerd op slides + oefeningen + studoforum examenvragen
% =============================================================================

\documentclass[a4paper,11pt]{article}
\usepackage{../../school-macros}
\usetikzlibrary{positioning, shapes.geometric, arrows.meta, fit, calc, backgrounds}
% SQL listing environment: avoids colupper+listings catcode conflict in codeblock
\lstdefinelanguage{MySQL}{
  keywords={CREATE,TABLE,PRIMARY,KEY,FOREIGN,REFERENCES,NOT,NULL,UNIQUE,
            DEFAULT,INT,VARCHAR,DATE,DATETIME,DECIMAL,TINYINT,TIME,FLOAT},
  sensitive=false,
  morecomment=[l]{--},
  morestring=[b]',
}
\lstnewenvironment{sqllisting}[1][]
{
  \lstset{
    language=MySQL,
    basicstyle=\ttfamily\small,
    keywordstyle=\bfseries\color{schoolBlue},
    commentstyle=\itshape\color{schoolGreen},
    breaklines=true, tabsize=4,
    frame=leftline, framerule=3pt, rulecolor=\color{schoolBlue!40},
    xleftmargin=15pt, numbers=left, numberstyle=\tiny\color{gray},
    showstringspaces=false, #1
  }
}
{}

% ERD TikZ stijlen
\tikzset{
  erd entity/.style={
    rectangle, draw=schoolBlue, thick, rounded corners=2pt,
    minimum width=2.8cm, minimum height=0.6cm,
    fill=schoolBlue!10, font=\bfseries\sffamily\small, align=center
  },
  erd attr/.style={
    rectangle, draw=schoolGray, thin,
    minimum width=2.8cm, minimum height=0.5cm,
    fill=gray!5, font=\ttfamily\footnotesize, align=left, inner sep=3pt
  },
  erd pk/.style={
    erd attr, draw=schoolOrange, fill=schoolOrange!8,
    font=\ttfamily\footnotesize\bfseries
  },
  erd fk/.style={
    erd attr, draw=schoolGreen, fill=schoolGreen!8,
    font=\ttfamily\footnotesize
  },
  erd rel/.style={
    draw=schoolGray!70, thick
  },
  one/.style={},
  many/.style={},
}

\newcommand{\erdbox}[3]{%
  % #1 = node name, #2 = entity name, #3 = attributes as comma list
  % usage: manually per diagram
}

\newcommand{\titel}{Statistiek \& Databeheer -- Database Design \& ERD}
\newcommand{\auteur}{KUL Industriële Ingenieurs}
\newcommand{\academiejaar}{Academiejaar 2025-2026}

\begin{document}

\kuleuventitle{\titel}{}{\auteur}{\academiejaar}

\begin{abstract}
Samenvatting van het databeheer-gedeelte met focus op examenvaardigheden:
een database ontwerpen vanuit een tekstbeschrijving via een ERD, en omzetten naar
SQL \texttt{CREATE TABLE}-statements in MySQL Workbench.
Het examen bestaat uit \textbf{±5 tabellen aanmaken} (DDL) en \textbf{1 SELECT-query} (DML).
\end{abstract}

\tableofcontents
\newpage

% =============================================================================
\section{Het Data Model}
% =============================================================================

\begin{conceptbox}[Proces- vs.\ Datamodel]
\begin{itemize}
  \item \textbf{Procesmodel} -- beschrijft \emph{hoe} een operatie uitgevoerd wordt.
  \item \textbf{Datamodel} -- beschrijft de \emph{structuur} van data: welke data, hoe opgeslagen, welke verbanden.
\end{itemize}
Het ERD is een \textbf{datamodel}. Het zegt niets over hoe data aangemaakt, gewijzigd of bevraagd wordt.
\end{conceptbox}

De aanpak voor elke database-oefening volgt altijd dezelfde stappen:

\begin{enumerate}
  \item Tekstbeschrijving lezen en \textbf{entiteiten} identificeren (zelfstandige naamwoorden).
  \item \textbf{Relaties} en \textbf{cardinaliteit} bepalen (werkwoorden/verbanden).
  \item \textbf{ERD} tekenen (mentaal of op papier -- niet op examen ingediend).
  \item \textbf{SQL DDL} schrijven: \texttt{CREATE TABLE} voor elke entiteit/koppeltabel.
\end{enumerate}

% =============================================================================
\section{ERD Bouwstenen}
% =============================================================================

\subsection{Entiteit}

\begin{theorie}[Entiteit (Entity)]
Een \concept{entiteit} is alles waarover we informatie willen opslaan = \textbf{zelfstandig naamwoord}.
\begin{itemize}
  \item Voorbeelden: \texttt{Student}, \texttt{Cursus}, \texttt{Boek}, \texttt{Branch}, \texttt{EscapeRoom}
  \item Wordt een \textbf{tabel} in de database.
  \item Een \textbf{instantie} is één concrete rij in die tabel (bv.\ de stad ``Leuven'').
\end{itemize}
\end{theorie}

\begin{waarschuwing}
Een entiteit = \emph{ding} waarover we data willen. Processen (boeken, betalen, reserveren) zijn \textbf{geen} entiteiten tenzij we er data van bijhouden.
\end{waarschuwing}

\subsubsection*{Zwakke entiteit (Weak Entity)}
Een entiteit die \textbf{geen eigen sleutelattribuut} kan hebben en altijd gekoppeld is aan een andere (sterke) entiteit via een 1:n-relatie.

\textit{Voorbeeld:} een \texttt{BookCopy} (exemplaar van een boek) bestaat niet zonder het boek zelf.

\subsubsection*{Associatieve entiteit}
Een entiteit die een \textbf{n:m-relatie} omzet naar twee 1:n-relaties. Ze bevat informatie over de relatie zelf.

\textit{Voorbeeld:} \texttt{Subscription} legt de relatie tussen \texttt{Student} en \texttt{Course} vast, en bevat ook \texttt{evaluation} en \texttt{startDate}.

\subsection{Attribuut}

\begin{theorie}[Attribuut]
Extra informatie over een entiteit. Wordt een \textbf{kolom} in de tabel.
\begin{itemize}
  \item \concept{Sleutelattribuut (Primary Key)} -- identificeert een instantie uniek.
  Bv.\ \texttt{studentNr}, \texttt{socialSecurityNumber}, \texttt{loginName}.
  \item \textbf{Normaal attribuut} -- beschrijvende info: \texttt{name}, \texttt{email}, \texttt{salary}.
\end{itemize}
\end{theorie}

\subsection{Relatie en Cardinaliteit}

\begin{theorie}[Relatie]
Een \concept{relatie} is een natuurlijke verbinding tussen twee entiteiten = \textbf{werkwoord}.
\begin{itemize}
  \item Voorbeelden: \emph{Student \underline{volgt} Cursus}, \emph{Branch \underline{heeft} EscapeRoom}
  \item Optioneel \(\leftrightarrow\) Verplicht (min.\ cardinaliteit = 0 of 1)
\end{itemize}
\end{theorie}

\begin{conceptbox}[Cardinaliteit -- de drie soorten]
\begin{center}
\begin{tblr}{
  colspec={Q[c,2cm] Q[l,4cm] Q[l,5.5cm]},
  row{1}={bg=schoolBlue, fg=white, font=\bfseries\sffamily},
  row{odd}={bg=gray!8},
}
  Type & Betekenis & Aanpak \\
  \textbf{1:1} & Eén instantie $\leftrightarrow$ één instantie & FK in één van de twee tabellen \\
  \textbf{1:n} & Eén $\leftrightarrow$ meerdere & FK in de \emph{veel}-kant tabel \\
  \textbf{n:m} & Meerdere $\leftrightarrow$ meerdere & Nieuwe \emph{koppeltabel} aanmaken \\
\end{tblr}
\end{center}
\end{conceptbox}

% =============================================================================
\section{Van ERD naar SQL-Tabellen}
% =============================================================================

\begin{waarschuwing}[Examen -- Dit is het kernstuk]
Op het examen krijg je een tekstbeschrijving en moet je de \texttt{CREATE TABLE}-statements schrijven.
Beheers deze drie regels uit je hoofd.
\end{waarschuwing}

\subsection{Regel 1 -- Entiteit $\rightarrow$ Tabel}

Elke entiteit wordt een tabel. Attributen worden kolommen. Het sleutelattribuut wordt \texttt{PRIMARY KEY}.

\begin{voorbeeld}[Student-entiteit]
\begin{sqllisting}
CREATE TABLE Student (
    studentNr   INT          PRIMARY KEY,
    name        VARCHAR(50)  NOT NULL,
    firstname   VARCHAR(50)  NOT NULL,
    birthdate   DATE         NOT NULL,
    email       VARCHAR(100) UNIQUE
);
\end{sqllisting}
\end{voorbeeld}

\subsection{Regel 2 -- 1:n Relatie $\rightarrow$ Foreign Key in de \emph{veel}-tabel}

De tabel aan de ``veel''-kant krijgt een \texttt{FOREIGN KEY} die verwijst naar de primaire sleutel van de ``één''-kant.

\begin{center}
\begin{tikzpicture}[node distance=0cm]
  % Branch box
  \node[erd entity] (branch) {Branch};
  \node[erd pk, below=of branch, anchor=north] (branchPK) {\underline{branchID} INT (PK)};
  \node[erd attr, below=of branchPK, anchor=north] (branchLoc) {location VARCHAR(100)};
  \node[erd attr, below=of branchLoc, anchor=north] (branchHrs) {openingHours VARCHAR(50)};
  \begin{scope}[on background layer]
    \node[draw=schoolBlue, rounded corners=2pt, thick, fill=schoolBlue!5,
          fit=(branch)(branchPK)(branchLoc)(branchHrs), inner sep=1pt] {};
  \end{scope}

  % EscapeRoom box
  \node[erd entity, right=3.5cm of branch] (room) {EscapeRoom};
  \node[erd pk, below=of room, anchor=north] (roomPK) {\underline{roomID} INT (PK)};
  \node[erd attr, below=of roomPK, anchor=north] (roomName) {roomName VARCHAR(100)};
  \node[erd attr, below=of roomName, anchor=north] (roomDiff) {difficulty VARCHAR(20)};
  \node[erd fk, below=of roomDiff, anchor=north] (roomFK) {branchID INT (FK)};
  \begin{scope}[on background layer]
    \node[draw=schoolBlue, rounded corners=2pt, thick, fill=schoolBlue!5,
          fit=(room)(roomPK)(roomName)(roomDiff)(roomFK), inner sep=1pt] {};
  \end{scope}

  % Arrow: 1 branch has n rooms
  \draw[erd rel, ->, >=Stealth] (roomFK.west) -- ++(-0.5,0) |- (branchPK.east)
    node[midway, above, font=\sffamily\scriptsize, text=schoolGray] {1:n};
\end{tikzpicture}
\end{center}

\begin{voorbeeld}[Branch heeft meerdere EscapeRooms]
\begin{sqllisting}
CREATE TABLE EscapeRoom (
    roomID      INT          PRIMARY KEY AUTO_INCREMENT,
    roomName    VARCHAR(100) NOT NULL,
    difficulty  VARCHAR(20)  NOT NULL,
    branchID    INT          NOT NULL,
    FOREIGN KEY (branchID) REFERENCES Branch(branchID)
);
\end{sqllisting}
\end{voorbeeld}

\subsection{Regel 3 -- n:m Relatie $\rightarrow$ Koppeltabel}

Bij een veel-op-veel-relatie maak je een \textbf{nieuwe tabel} aan die de primaire sleutels van beide tabellen als foreign keys bevat. Die tabel kan ook extra attributen hebben (bv.\ datum, prijs).

\begin{center}
\begin{tikzpicture}[node distance=0cm]
  % Customer
  \node[erd entity] (cust) {Customer};
  \node[erd pk, below=of cust, anchor=north] (custPK) {\underline{customerID} INT (PK)};
  \node[erd attr, below=of custPK, anchor=north] (custName) {fullName VARCHAR(100)};
  \node[erd attr, below=of custName, anchor=north] (custAge) {birthdate DATE};
  \begin{scope}[on background layer]
    \node[draw=schoolBlue, rounded corners=2pt, thick, fill=schoolBlue!5,
          fit=(cust)(custPK)(custName)(custAge), inner sep=1pt] {};
  \end{scope}

  % CustomerParty (koppeltabel)
  \node[erd entity, right=3cm of cust, fill=schoolOrange!20, draw=schoolOrange] (cp) {CustomerParty};
  \node[erd fk, below=of cp, anchor=north, draw=schoolOrange, fill=schoolOrange!8] (cpCust) {customerID INT (FK/PK)};
  \node[erd fk, below=of cpCust, anchor=north, draw=schoolOrange, fill=schoolOrange!8] (cpParty) {partyID INT (FK/PK)};
  \begin{scope}[on background layer]
    \node[draw=schoolOrange, rounded corners=2pt, thick, fill=schoolOrange!5,
          fit=(cp)(cpCust)(cpParty), inner sep=1pt] {};
  \end{scope}

  % Party
  \node[erd entity, right=3cm of cp] (party) {Party};
  \node[erd pk, below=of party, anchor=north] (partyPK) {\underline{partyID} INT (PK)};
  \begin{scope}[on background layer]
    \node[draw=schoolBlue, rounded corners=2pt, thick, fill=schoolBlue!5,
          fit=(party)(partyPK), inner sep=1pt] {};
  \end{scope}

  \draw[erd rel, ->, >=Stealth] (cpCust.west) -- ++(-0.5,0) |- (custPK.east);
  \draw[erd rel, ->, >=Stealth] (cpParty.east) -- ++(0.5,0) |- (partyPK.west);
\end{tikzpicture}
\end{center}

\begin{voorbeeld}[Customer $\leftrightarrow$ Party via koppeltabel]
\begin{sqllisting}
CREATE TABLE CustomerParty (
    customerID  INT  NOT NULL,
    partyID     INT  NOT NULL,
    PRIMARY KEY (customerID, partyID),
    FOREIGN KEY (customerID) REFERENCES Customer(customerID),
    FOREIGN KEY (partyID)    REFERENCES Party(partyID)
);
\end{sqllisting}
\end{voorbeeld}

\subsection{Overzichtstabel: ERD-element $\rightarrow$ SQL}

\begin{center}
\begin{tblr}{
  colspec={Q[l,4cm] Q[l,8cm]},
  row{1}={bg=schoolBlue, fg=white, font=\bfseries\sffamily},
  row{odd}={bg=gray!8},
}
  ERD-element & SQL-equivalent \\
  Entiteit & \texttt{CREATE TABLE Naam (\ldots)} \\
  Attribuut & Kolom in de tabel \\
  Sleutelattribuut & \texttt{PRIMARY KEY} \\
  1:1 of 1:n relatie & \texttt{FOREIGN KEY} in de ,,veel''-kant tabel \\
  n:m relatie & Nieuwe koppeltabel met 2 \texttt{FOREIGN KEY}s \\
  Associatieve entiteit & Koppeltabel + extra kolommen \\
  Zwakke entiteit & Tabel + verplichte \texttt{FOREIGN KEY} naar sterke entiteit \\
\end{tblr}
\end{center}

% =============================================================================
\section{SQL DDL -- CREATE TABLE Syntax}
% =============================================================================

\subsection{Datatypes}

\begin{center}
\begin{tblr}{
  colspec={Q[l,3.5cm] Q[l,8.5cm]},
  row{1}={bg=schoolBlue, fg=white, font=\bfseries\sffamily},
  row{odd}={bg=gray!8},
}
  Datatype & Gebruik \\
  \texttt{INT} & Geheel getal (ook \texttt{INTEGER}) \\
  \texttt{TINYINT(1)} & Boolean (0/1) -- MySQL heeft geen echte bool \\
  \texttt{VARCHAR(n)} & Tekst met max.\ \texttt{n} tekens \\
  \texttt{DATE} & Datum: \texttt{YYYY-MM-DD} \\
  \texttt{TIME} & Tijd: \texttt{HH:MM:SS} \\
  \texttt{DATETIME} & Datum + tijd samen \\
  \texttt{DECIMAL(p,s)} & Decimaal getal: \texttt{p} cijfers totaal, \texttt{s} na komma \\
  \texttt{FLOAT} / \texttt{DOUBLE} & Vlottend getal (minder precies dan DECIMAL) \\
\end{tblr}
\end{center}

\subsection{Constraints (beperkingen)}

\begin{center}
\begin{tblr}{
  colspec={Q[l,3.5cm] Q[l,8.5cm]},
  row{1}={bg=schoolBlue, fg=white, font=\bfseries\sffamily},
  row{odd}={bg=gray!8},
}
  Constraint & Betekenis \\
  \texttt{PRIMARY KEY} & Uniek + niet null; identificeert rij \\
  \texttt{NOT NULL} & Verplicht in te vullen \\
  \texttt{UNIQUE} & Waarde moet uniek zijn in de kolom \\
  \texttt{DEFAULT x} & Standaardwaarde als niets opgegeven \\
  \texttt{AUTO\_INCREMENT} & Automatisch verhogend getal (voor PK) \\
  \texttt{FOREIGN KEY \ldots REFERENCES} & Verwijst naar PK van andere tabel \\
  \texttt{CHECK (voorwaarde)} & Waarde moet aan voorwaarde voldoen \\
\end{tblr}
\end{center}

\subsection{Volledige Syntaxtemplate}

\begin{sqllisting}
CREATE TABLE TabelNaam (
    -- Kolommen
    kolomNaam1   DATATYPE  [NOT NULL] [UNIQUE] [DEFAULT waarde],
    kolomNaam2   DATATYPE  [NOT NULL],
    -- Primary key (ofwel inline, ofwel hier)
    PRIMARY KEY (kolomNaam1),
    -- Foreign keys (altijd aan het einde)
    FOREIGN KEY (kolomNaamFK) REFERENCES AndereTabel(anderePK)
);
\end{sqllisting}

\begin{waarschuwing}[Volgorde in CREATE TABLE]
\begin{enumerate}
  \item Eerst alle gewone kolommen
  \item Dan \texttt{PRIMARY KEY}
  \item Dan \texttt{FOREIGN KEY} (altijd als laatste)
  \item Komma na elke regel \textbf{behalve de laatste}
\end{enumerate}
\end{waarschuwing}

\subsection{Foreign Key -- REFERENCES}

\begin{theorie}[Foreign Key regel]
Als je een waarde invult voor een foreign key, \textbf{moet} de gerefereerde rij bestaan in de andere tabel.
Elke waarde van een FK moet overeenkomen met de PK van een rij in de gerefereerde tabel.
\end{theorie}

\begin{voorbeeld}[FK syntax met ON DELETE]
\begin{sqllisting}
FOREIGN KEY (branchID)
    REFERENCES Branch(branchID)
    ON DELETE CASCADE    -- als Branch verwijderd wordt, ook EscapeRoom verwijderen
    ON UPDATE CASCADE    -- als branchID wijzigt, automatisch updaten
\end{sqllisting}
\end{voorbeeld}

% =============================================================================
\section{Normalisatie}
% =============================================================================

\begin{conceptbox}[Normalisatie -- 3 regels]
Normalisatie maakt een datamodel \textbf{flexibel en betrouwbaar}. Er zijn 3 regels (normaalvormen):
\begin{enumerate}
  \item \textbf{1NV: Geen herhalende groepen} -- Als een entiteit meerdere waarden voor een attribuut heeft (bv.\ meerdere telefoonnummers), maak er een aparte entiteit van via 1:n.
  \item \textbf{2NV: Volledige functionele afhankelijkheid} -- Alle attributen moeten afhankelijk zijn van de \emph{volledige} primaire sleutel (relevant bij samengestelde PK).
  \item \textbf{3NV: Geen transitieve afhankelijkheid} -- Attributen mogen niet afhankelijk zijn van een niet-sleutelattribuut. Bv.\ als \texttt{zipCode} de \texttt{city} bepaalt, maak dan een aparte \texttt{City}-tabel.
\end{enumerate}
\end{conceptbox}

\begin{waarschuwing}
Meer tabellen door normalisatie $\Rightarrow$ meer overhead bij queries (JOINs nodig), maar betere datamodellering.
\end{waarschuwing}

% =============================================================================
\section{Uitgewerkte Oefening -- Escape Room}
% =============================================================================

\begin{oefening}[Escape Room: Bolting from Belgium]
Het escape room bedrijf ``Bolting from Belgium'' heeft meerdere \textbf{vestigingen (branches)} in België, elk met eigen openingstijden en locatie. Elke vestiging heeft meerdere \textbf{escape rooms} met een naam en moeilijkheidsgraad. Klanten bezoeken een escape room in \textbf{groepen (parties)}. \textbf{Klanten} worden geregistreerd met volledige naam en geboortedatum. Een klant kan tot meerdere parties behoren. Wanneer een party een kamer boekt, worden datum, begintijd en prijs opgeslagen.
\end{oefening}

\subsection{Stap 1 -- Entiteiten identificeren}

Zelfstandige naamwoorden uit de tekst:
\textbf{Branch} · \textbf{EscapeRoom} · \textbf{Customer} · \textbf{Party} · \textbf{Booking}

Relaties:
\begin{itemize}
  \item Branch \emph{heeft} EscapeRoom: \textbf{1:n} (één vestiging, meerdere kamers)
  \item Customer \emph{behoort tot} Party: \textbf{n:m} (klant in meerdere parties, party met meerdere klanten)
  \item Party \emph{boekt} EscapeRoom: \textbf{n:m met extra data} (datum, begintijd, prijs) $\Rightarrow$ \textbf{Booking} als associatieve entiteit
\end{itemize}

\subsection{Stap 2 -- ERD Schets}

\begin{center}
\begin{tikzpicture}[node distance=0.5cm and 3.2cm, every node/.style={font=\sffamily\small}]

  % Branch
  \node[erd entity] (branch) {Branch};
  \node[erd pk, below=0cm of branch] (bPK) {\underline{branchID}};
  \node[erd attr, below=0cm of bPK] (bLoc) {location};
  \node[erd attr, below=0cm of bLoc] (bHrs) {openingHours};
  \begin{scope}[on background layer]
    \node[draw=schoolBlue, rounded corners=2pt, thick, fill=schoolBlue!5,
          fit=(branch)(bPK)(bLoc)(bHrs), inner sep=2pt] {};
  \end{scope}

  % EscapeRoom
  \node[erd entity, right=3.5cm of branch] (room) {EscapeRoom};
  \node[erd pk, below=0cm of room] (rPK) {\underline{roomID}};
  \node[erd attr, below=0cm of rPK] (rName) {roomName};
  \node[erd attr, below=0cm of rName] (rDiff) {difficulty};
  \node[erd fk, below=0cm of rDiff] (rFK) {branchID (FK)};
  \begin{scope}[on background layer]
    \node[draw=schoolBlue, rounded corners=2pt, thick, fill=schoolBlue!5,
          fit=(room)(rPK)(rName)(rDiff)(rFK), inner sep=2pt] {};
  \end{scope}

  % Booking (associatief)
  \node[erd entity, right=3.5cm of room, fill=schoolOrange!20, draw=schoolOrange] (booking) {Booking};
  \node[erd fk, below=0cm of booking, draw=schoolOrange, fill=schoolOrange!8] (bkParty) {partyID (FK/PK)};
  \node[erd fk, below=0cm of bkParty, draw=schoolOrange, fill=schoolOrange!8] (bkRoom) {roomID (FK/PK)};
  \node[erd attr, below=0cm of bkRoom] (bkDate) {bookingDate};
  \node[erd attr, below=0cm of bkDate] (bkTime) {startTime};
  \node[erd attr, below=0cm of bkTime] (bkPrice) {price};
  \begin{scope}[on background layer]
    \node[draw=schoolOrange, rounded corners=2pt, thick, fill=schoolOrange!5,
          fit=(booking)(bkParty)(bkRoom)(bkDate)(bkTime)(bkPrice), inner sep=2pt] {};
  \end{scope}

  % Party
  \node[erd entity, below=2.8cm of booking] (party) {Party};
  \node[erd pk, below=0cm of party] (pPK) {\underline{partyID}};
  \begin{scope}[on background layer]
    \node[draw=schoolBlue, rounded corners=2pt, thick, fill=schoolBlue!5,
          fit=(party)(pPK), inner sep=2pt] {};
  \end{scope}

  % CustomerParty (koppeltabel n:m)
  \node[erd entity, below=2.8cm of room, fill=schoolTeal!20, draw=schoolTeal] (cp) {CustomerParty};
  \node[erd fk, below=0cm of cp, draw=schoolTeal, fill=schoolTeal!8] (cpC) {customerID (FK/PK)};
  \node[erd fk, below=0cm of cpC, draw=schoolTeal, fill=schoolTeal!8] (cpP) {partyID (FK/PK)};
  \begin{scope}[on background layer]
    \node[draw=schoolTeal, rounded corners=2pt, thick, fill=schoolTeal!5,
          fit=(cp)(cpC)(cpP), inner sep=2pt] {};
  \end{scope}

  % Customer
  \node[erd entity, below=2.8cm of branch] (cust) {Customer};
  \node[erd pk, below=0cm of cust] (cPK) {\underline{customerID}};
  \node[erd attr, below=0cm of cPK] (cName) {fullName};
  \node[erd attr, below=0cm of cName] (cBirth) {birthdate};
  \begin{scope}[on background layer]
    \node[draw=schoolBlue, rounded corners=2pt, thick, fill=schoolBlue!5,
          fit=(cust)(cPK)(cName)(cBirth), inner sep=2pt] {};
  \end{scope}

  % Relaties
  \draw[erd rel, ->, >=Stealth] (rFK.west) -- (bLoc.east |- rFK.west)
    node[midway, above, font=\scriptsize\sffamily, text=schoolGray] {1:n};
  \draw[erd rel, ->, >=Stealth] (bkRoom.west) -- (rPK.east |- bkRoom.west)
    node[midway, above, font=\scriptsize\sffamily, text=schoolGray] {};
  \draw[erd rel, ->, >=Stealth] (bkParty.south) -- (party.north |- bkParty.south)
    node[midway, right, font=\scriptsize\sffamily, text=schoolGray] {};
  \draw[erd rel, ->, >=Stealth] (cpC.west) -- (cPK.east |- cpC.west);
  \draw[erd rel, ->, >=Stealth] (cpP.east) -- (pPK.west |- cpP.east);

\end{tikzpicture}
\end{center}

\subsection{Stap 3 -- SQL Oplossing}

\begin{sqllisting}
-- Tabel 1: Branch (geen FK's nodig)
CREATE TABLE Branch (
    branchID      INT          PRIMARY KEY AUTO_INCREMENT,
    location      VARCHAR(100) NOT NULL,
    openingHours  VARCHAR(100) NOT NULL
);

-- Tabel 2: EscapeRoom (FK naar Branch)
CREATE TABLE EscapeRoom (
    roomID      INT          PRIMARY KEY AUTO_INCREMENT,
    roomName    VARCHAR(100) NOT NULL,
    difficulty  VARCHAR(20)  NOT NULL,
    branchID    INT          NOT NULL,
    FOREIGN KEY (branchID) REFERENCES Branch(branchID)
);

-- Tabel 3: Customer
CREATE TABLE Customer (
    customerID  INT          PRIMARY KEY AUTO_INCREMENT,
    fullName    VARCHAR(100) NOT NULL,
    birthdate   DATE         NOT NULL
);

-- Tabel 4: Party
CREATE TABLE Party (
    partyID  INT  PRIMARY KEY AUTO_INCREMENT
);

-- Koppeltabel 5: CustomerParty  (n:m Customer <-> Party)
CREATE TABLE CustomerParty (
    customerID  INT  NOT NULL,
    partyID     INT  NOT NULL,
    PRIMARY KEY (customerID, partyID),
    FOREIGN KEY (customerID) REFERENCES Customer(customerID),
    FOREIGN KEY (partyID)    REFERENCES Party(partyID)
);

-- Associatieve tabel 6: Booking  (Party boekt EscapeRoom + extra data)
CREATE TABLE Booking (
    partyID      INT      NOT NULL,
    roomID       INT      NOT NULL,
    bookingDate  DATE     NOT NULL,
    startTime    TIME     NOT NULL,
    price        DECIMAL(8,2) NOT NULL,
    PRIMARY KEY (partyID, roomID, bookingDate, startTime),
    FOREIGN KEY (partyID) REFERENCES Party(partyID),
    FOREIGN KEY (roomID)  REFERENCES EscapeRoom(roomID)
);
\end{sqllisting}

\begin{waarschuwing}[Let op: volgorde van aanmaken]
Tabellen met \texttt{FOREIGN KEY}s moeten aangemaakt worden \emph{nadat} de gerefereerde tabel al bestaat.
Volgorde hier: \texttt{Branch} $\to$ \texttt{EscapeRoom} $\to$ \texttt{Customer} $\to$ \texttt{Party} $\to$ \texttt{CustomerParty} $\to$ \texttt{Booking}.
\end{waarschuwing}

% =============================================================================
\section{Uitgewerkte Oefening -- Basketbalclub}
% =============================================================================

\begin{oefening}[Basketbalclub]
Een basketbalclub wil spelersinformatie bijhouden. Een \textbf{club} heeft naam, adres, voorzitter en registratienummer. Een club heeft meerdere \textbf{teams} per categorie (jeugd, junioren\ldots) met een naam en trainer. Een team heeft \textbf{spelers} met naam, geboortedatum en ranking. Per \textbf{match} (naam, tijdstip, eindstand) -- gespeeld tegen een team van een andere club -- worden per speler de gespeelde tijd, individuele score en aantal 3-punters bijgehouden.
\end{oefening}

\subsection{Entiteiten en Relaties}

\begin{itemize}
  \item \textbf{Club} -- naam, adres, voorzitter, registratienummer
  \item \textbf{Team} -- naam, categorie, trainer; \textbf{FK naar Club} (1:n)
  \item \textbf{Player} -- naam, geboortedatum, ranking; \textbf{FK naar Team} (1:n)
  \item \textbf{Match} -- naam, tijdstip, eindstand; gespeeld door een Team tegen een andere Club
  \item \textbf{PlayerMatchStats} -- koppeltabel Player $\leftrightarrow$ Match met extra data (gespeelde tijd, score, 3-punters)
\end{itemize}

\subsection{SQL Oplossing}

\begin{sqllisting}
CREATE TABLE Club (
    clubID       INT          PRIMARY KEY AUTO_INCREMENT,
    clubName     VARCHAR(100) NOT NULL,
    address      VARCHAR(200) NOT NULL,
    president    VARCHAR(100) NOT NULL,
    leagueNr     VARCHAR(50)  UNIQUE NOT NULL
);

CREATE TABLE Team (
    teamID    INT          PRIMARY KEY AUTO_INCREMENT,
    teamName  VARCHAR(100) NOT NULL,
    category  VARCHAR(50)  NOT NULL,
    coach     VARCHAR(100) NOT NULL,
    clubID    INT          NOT NULL,
    FOREIGN KEY (clubID) REFERENCES Club(clubID)
);

CREATE TABLE Player (
    playerID   INT          PRIMARY KEY AUTO_INCREMENT,
    playerName VARCHAR(100) NOT NULL,
    birthdate  DATE         NOT NULL,
    ranking    INT          NOT NULL,
    teamID     INT          NOT NULL,
    FOREIGN KEY (teamID) REFERENCES Team(teamID)
);

CREATE TABLE Match (
    matchID     INT          PRIMARY KEY AUTO_INCREMENT,
    matchName   VARCHAR(100) NOT NULL,
    matchTime   DATETIME     NOT NULL,
    finalScore  VARCHAR(20),
    homeTeamID  INT          NOT NULL,        -- het team van onze club
    opponentID  INT          NOT NULL,        -- team van andere club
    FOREIGN KEY (homeTeamID) REFERENCES Team(teamID),
    FOREIGN KEY (opponentID) REFERENCES Team(teamID)
);

-- Associatieve tabel: spelersstats per match
CREATE TABLE PlayerMatchStats (
    playerID      INT  NOT NULL,
    matchID       INT  NOT NULL,
    minutesPlayed INT,
    individualScore INT,
    threePointers INT,
    PRIMARY KEY (playerID, matchID),
    FOREIGN KEY (playerID) REFERENCES Player(playerID),
    FOREIGN KEY (matchID)  REFERENCES Match(matchID)
);
\end{sqllisting}

% =============================================================================
\section{Snelreferentie -- Checklist voor het Examen}
% =============================================================================

\begin{conceptbox}[Stap-voor-stap aanpak op het examen]
\begin{enumerate}
  \item \textbf{Lees de tekst} en onderstreep alle zelfstandige naamwoorden $\Rightarrow$ kandidaat-entiteiten.
  \item \textbf{Zoek relaties} (werkwoorden) en bepaal de cardinaliteit:
    \begin{itemize}
      \item Kan één X meerdere Y's hebben? $\Rightarrow$ 1:n (FK in Y)
      \item Kan ook Y meerdere X's hebben? $\Rightarrow$ n:m (koppeltabel)
      \item Is er extra data bij de relatie (datum, prijs\ldots)? $\Rightarrow$ associatieve entiteit/tabel
    \end{itemize}
  \item \textbf{Bepaal de volgorde} van aanmaken: tabellen zonder FK's eerst.
  \item \textbf{Schrijf de SQL}: voor elke entiteit een \texttt{CREATE TABLE}.
  \item \textbf{Controleer}: elke FK verwijst naar een bestaande PK? Komma's correct?
\end{enumerate}
\end{conceptbox}

\begin{waarschuwing}[Veelgemaakte fouten op het examen]
\begin{itemize}
  \item FK vergeten bij 1:n-relatie (staat in de \emph{veel}-tabel, niet in de één-tabel).
  \item Geen koppeltabel aanmaken bij n:m (dit is de meest voorkomende fout).
  \item Komma na de \emph{laatste} kolom/constraint (syntaxfout).
  \item Tabellen aanmaken in verkeerde volgorde (FK verwijst naar nog niet bestaande tabel).
  \item \texttt{REFERENCES TabelNaam(PKkolom)} -- moet de exacte kolomnaam van de PK zijn.
\end{itemize}
\end{waarschuwing}

\begin{center}
\begin{tblr}{
  colspec={Q[l,3.5cm] Q[l,8.5cm]},
  row{1}={bg=schoolBlue, fg=white, font=\bfseries\sffamily},
  row{odd}={bg=gray!8},
}
  Situatie & Oplossing \\
  1 entiteit & 1 \texttt{CREATE TABLE} met \texttt{PRIMARY KEY} \\
  1:n relatie & FK in de \emph{veel}-tabel \texttt{REFERENCES} één-tabel \\
  n:m relatie & Koppeltabel met 2 FK's en composite \texttt{PRIMARY KEY} \\
  n:m + extra data & Associatieve tabel = koppeltabel + extra kolommen \\
  Verplicht veld & \texttt{NOT NULL} \\
  Uniek veld & \texttt{UNIQUE} \\
  Auto-gegenereerde ID & \texttt{AUTO\_INCREMENT} bij \texttt{INT PRIMARY KEY} \\
  Tekst & \texttt{VARCHAR(n)} -- kies n iets groter dan nodig \\
  Datum & \texttt{DATE}, Datum+Tijd: \texttt{DATETIME} \\
  Prijs/bedrag & \texttt{DECIMAL(8,2)} \\
\end{tblr}
\end{center}

\end{document}
